\section*{Chapitre \ref{ch:hf:crypto-high-frequency-trading-on-centralised-venues} --- Trading à haute fréquence sur plateformes centralisées de crypto-actifs}

\begin{solution}{exo:hf:crypto-high-frequency-trading-on-centralised-venues:1}
$2+0.5-(0.5+1)=1.0$ point de base par exécution.
\end{solution}

\begin{solution}{exo:hf:crypto-high-frequency-trading-on-centralised-venues:2}
$200\,000/86\,400=2.31$ ordres par seconde, 1.16 recotations~; pour cinq instruments, 0.23 chacun.
\end{solution}

\begin{solution}{exo:hf:crypto-high-frequency-trading-on-centralised-venues:3}
Les positions longues paient les positions courtes lorsque le taux est positif~: la position courte reçoit $150\,000\times0.0001=\$15$.
\end{solution}

\begin{solution}{exo:hf:crypto-high-frequency-trading-on-centralised-venues:4}
Dès la première seconde, l'ordre à cours limité s'exécute à de faibles décotes, et la position subit ensuite le reste de la baisse avant de remonter~; attendre jusqu'à 40
points de base permet d'acheter à des décotes plus profondes et de capter toute la remontée à partir de là~: 20\,800~\$ contre 5\,200~\$.
\end{solution}

\begin{solution}{exo:hf:crypto-high-frequency-trading-on-centralised-venues:5}
Entre deux recotations, le \omterm{def:hf:fair-value:fair}{prix juste} évolue~; lorsqu'il varie de plus que la demi-fourchette augmentée du seuil des arbitragistes, ces derniers saisissent la cotation périmée à un
prix déjà moins favorable que le \omterm{def:hf:fair-value:fair}{prix juste}. La fourchette ne rapporte que sur les exécutions à des cotations fraîches.
\end{solution}

\begin{solution}{exo:hf:crypto-high-frequency-trading-on-centralised-venues:6}
Les contrôles des capitaux et la difficulté à transférer des fonds d'un pays à l'autre, qui empêchaient les arbitragistes de refermer les écarts entre pays.
\end{solution}

\begin{solution}{exo:hf:crypto-high-frequency-trading-on-centralised-venues:7}
Le seuil de 60 points de base rapporte 17\,800~\$ grâce à la cascade~; la limite de 2~millions de dollars, active dès la première seconde, perd 600~\$. Le moment compte davantage que
la capacité~: la petite position initiale est complète avant que la décote vaille la peine d'être prise.
\end{solution}

\begin{solution}{exo:hf:crypto-high-frequency-trading-on-centralised-venues:8}
Le plafond quotidien d'ordres d'un compte autorise environ 0.06 recotation par seconde et par instrument, pour vingt instruments~; recoter chacun à chaque seconde serait refusé, et au
rythme autorisé, les cotations seraient périmées et saisies par les arbitragistes. Le backtest doit respecter les limites de la plateforme.
\end{solution}

\begin{solution}{pb:hf:crypto-high-frequency-trading-on-centralised-venues:1}
\begin{enumerate}
\item Voir la \cref{def:hf:crypto-high-frequency-trading-on-centralised-venues:xmm} et la \cref{def:hf:crypto-high-frequency-trading-on-centralised-venues:hedge}.
\item Des opportunités d'arbitrage importantes et récurrentes entre bourses, plus grandes entre pays qu'à l'intérieur d'un même pays, maintenues ouvertes par les contrôles des capitaux.
\item 6\,000 de poids de requête par minute, 100 ordres toutes les 10 secondes, 200\,000 par jour~; taux de financement à 00:00, 08:00 et 16:00 UTC~; 19~milliards de dollars liquidés le 10 octobre
2025.
\item Un crypto-actif à 60\,000~\$, une volatilité quotidienne de 3~\% , des clients toutes les demi-secondes, des arbitragistes sur cotations périmées, 2 points de base de chaque côté, une couverture au comptant,
un financement toutes les huit heures, une cascade le quatrième jour.
\item La base entre le contrat perpétuel et le comptant, et le financement.
\item Voir la \cref{def:hf:crypto-high-frequency-trading-on-centralised-venues:absorb}.
\item 600~\$, 5\,200~\$, 20\,800~\$ et 17\,800~\$.
\item Une cascade plus profonde, la dislocation devenant le prix, l'auto-deleveraging (ADL), la solvabilité de la plateforme.
\item Le plafond quotidien~: 1.16 recotations par seconde pour un instrument.
\item 275\,400~\$, 270\,800~\$, 173\,000~\$ et $-\$265\,900$.
\item Les deux côtés, dans une fourchette maximale et avec une taille minimale~; les limites d'inventaire imposent de retirer un côté.
\item Sur les cotations les plus périmées par rapport au \omterm{def:hf:fair-value:fair}{prix juste}, et sur moins d'instruments.
\item Fourchette 574\,900~\$, couverture $-\$464\,300$, remises 143\,700~\$, financement 300~\$, cascade 20\,800~\$~: 275\,400~\$~; un budget d'une recotation toutes les quatre
secondes en abandonne 102\,400~\$, toutes les dix secondes la totalité et davantage.
\item La fourchette est la plus large~; le coût de la couverture et le budget sont les plus fragiles.
\item Un second \omterm{def:hf:fair-value:fair}{prix juste} et une seconde \omterm{def:hf:crypto-high-frequency-trading-on-centralised-venues:hedge}{plateforme de couverture}, ce qui réduit le glissement et le risque de plateforme.
\item L'\omterm{def:hf:crypto-high-frequency-trading-on-centralised-venues:absorb}{absorption des liquidations}, qui détient des positions sur les plateformes pendant leurs pires moments.
\item Porter le carry du côté qui reçoit à travers les horodatages lorsque la base le permet~; solder le côté qui paie avant ces horodatages.
\item Disponibilité et profondeur, tels que le cahier des charges les définit.
\item Comparer la trajectoire de base du modèle et le flux forcé avec les relevés des plateformes de ce jour-là.
\item L'immédiateté sur une plateforme, évaluée à partir d'une autre.
\end{enumerate}
\end{solution}

\begin{solution}{iq:hf:crypto-high-frequency-trading-on-centralised-venues:1}
Les frais sur les deux plateformes, la possibilité de transférer les fonds et les crypto-actifs (statut des retraits, délai de transfert), la profondeur derrière les cours, et si le cours d'une
plateforme est périmé ou suspendu.
\iqlookfor{transférabilité et profondeur exploitable.}
\end{solution}

\begin{solution}{iq:hf:crypto-high-frequency-trading-on-centralised-venues:2}
Un limiteur par plateforme et par compte, reproduisant chaque limite sous forme de fenêtre~; une file de priorité des requêtes selon leur valeur~; les couvertures avant les recotations~; une
marge de sécurité sous la limite~; réagir au code 429 en suspendant les envois.
\iqlookfor{priorités et marges.}
\end{solution}

\begin{solution}{iq:hf:crypto-high-frequency-trading-on-centralised-venues:3}
Le taux est une moyenne temporelle de l'indice de prime sur l'intervalle, à laquelle s'ajoute un terme d'intérêt borné~; l'essentiel en est connu avant l'horodatage, à mesure que la
moyenne se constitue.
\iqlookfor{la moyenne de la formule.}
\end{solution}

\begin{solution}{iq:hf:crypto-high-frequency-trading-on-centralised-venues:4}
Les crypto-actifs et les positions couvertes ailleurs contre eux~; des limites de soldes par plateforme selon le risque de plateforme, et des couvertures qui ne dépendent pas de
cette plateforme.
\iqlookfor{risque de contrepartie des plateformes.}
\end{solution}

\begin{solution}{iq:hf:crypto-high-frequency-trading-on-centralised-venues:5}
Seulement lorsque la base est profonde et que le flux forcé est visible, dans une capacité fixée par la taille probable de la remontée~; couvrir sur la plateforme au comptant la plus profonde.
\iqlookfor{moment et capacité.}
\end{solution}

\begin{solution}{iq:hf:crypto-high-frequency-trading-on-centralised-venues:6}
La variation du prix sur $\tau$ a pour écart-type $\sigma\sqrt\tau$~; la probabilité qu'elle dépasse $h$ en valeur absolue est $2(1-\Phi(h/(\sigma\sqrt\tau)))$,
si bien que des cotations périmées apparaissent à peu près à ce taux à chaque intervalle de recotation.
\iqlookfor{diffusion sur l'intervalle de recotation.}
\end{solution}
